\section{Ramification recovered from multiplication failure}
\label{sec:v122-introduction}
The reduced support of a multiplication-failure scheme need not determine
its embedded geometry. We show that a fixed rational reduced divisor can
underlie failure schemes from which one recovers a smooth genus-one curve
or a K3 surface. The recovering invariant is intrinsic to the nonreduced
scheme: it is the annihilator of its nilradical, followed by birational
cancellation. Thus neither an ambient framing nor a choice of embedded
primary component is needed for the recovery statement.

Here is the construction. Let $e\in\{3,4\}$, put $p=\binom e2$, and
let $V$ be an $e$-dimensional complex vector space. An $e$-plane
$R\subset\Sym^2V$ of quadratic relations defines
\[
 S_R=\Sym^2V/R,\qquad
 B_R=\C\oplus V\oplus S_R,
\]
where multiplication in degree two is the quotient map and all products
of degree at least three vanish. We specify a nonempty invariant open
$G_e^\circ\subset\Gr(e,\Sym^2V)$ in
Proposition~\ref{prop:admissible-ramification}. Its points are
basepoint-free systems of quadrics on $\Pj(V^*)$ with smooth ramification
hypersurface $Y_R$. The associated finite morphism
\[
 \phi_R:\Pj(V^*)\longrightarrow\Pj(R^*)
\]
has degree $2^{e-1}$, and $Y_R$ has degree $e$. Thus $Y_R$ is a
smooth plane cubic for $e=3$ and a smooth quartic K3 surface for $e=4$.
The class is an open set of a Grassmannian, not a finite list of algebra
orbits: the orbit-dimension deficits are at least one and nine,
respectively.

On $X_R=\Gr(e,V\oplus S_R)$ let $U_R$ be the open where $K\to V$
has rank at least $e-1$, and let $\Delta_R\subset U_R$ be its
determinant divisor. The inverse image $E_R$ of $Y_R$ under
$K\mapsto\operatorname{im}(K\to V)$ is smooth and integral of
codimension two. For every $m\ge2$, the failure scheme of
$\Sym^m(\C\oplus K)\to B_R$ on \emph{all} of $U_R$ satisfies
\begin{equation}\label{eq:v122-intro-primary}
 \mathcal I_{D_R}=\mathcal I_{\Delta_R}\mathcal I_{E_R}
       =\mathcal I_{\Delta_R}\cap\mathcal I_{E_R}^2.
\end{equation}
These are its two associated supports. More generally, for every $n\ge1$,
\begin{equation}\label{eq:v122-intro-powers}
 \mathcal I_{D_R}^{\,n}
       =\mathcal I_{\Delta_R}^{\,n}\cap\mathcal I_{E_R}^{\,2n}.
\end{equation}
Both decompositions are irredundant. Theorem~\ref{thm:ramification-primary}
proves these equalities directly for each fixed tensor, and
Theorem~\ref{thm:primary-powers-intrinsic} determines the embedded torsion,
its annihilator and its generic length for every $n$. These are powers
of an ideal, not additional symmetric multiplication degrees.

Writing $\mathcal N_R$ for the nilradical of $\mathcal O_{D_R}$, we have
\[
 \mathcal N_R^2=0,\qquad
 \mathcal N_R\simeq\mathcal O_{E_R}(-\Delta_R),\qquad
 E_R=V\!\left(\operatorname{Ann}_{\mathcal O_{D_R}}\mathcal N_R\right).
\]
The support $E_R$ is a vector bundle over $Y_R\times\Pj^{p-1}$.
It follows that an abstract isomorphism $D_R\simeq D_{R'}$ forces
$Y_R\simeq Y_{R'}$ (Theorem~\ref{thm:intrinsic-ramification-recovery}).
The reduced divisors, in contrast, are all isomorphic after identifying
the underlying vector spaces and have rational smooth projective
compactifications. For the explicit elliptic family retained below,
distinct $j$-invariants therefore give nonisomorphic \emph{full failure
schemes}, not merely nonisomorphic input algebras.

The link between the two geometries is an equality of cokernel sheaves.
For $H=\ker\alpha$, contraction gives
\[
 \operatorname{coker}\bigl(\Sym^2H\longrightarrow S_R\bigr)
 \simeq
 \operatorname{coker}\bigl(R\xrightarrow{\,c_\alpha\,}V\bigr).
\]
On projective space the target on the right carries its natural
$\mathcal O(1)$ twist. In a basis of quadrics the right-hand map is their
Jacobian matrix; its determinant defines the ramification of $\phi_R$.
This cokernel identity, the full multiplication Fitting ideal, and the
intrinsic annihilator are the successive steps of the recovery theorem.
They do not follow from the reduced determinant alone.

The primary theorem is also relative over $G_e^\circ$: the failure
schemes and all the specified ideal-power thickenings form flat families
(Proposition~\ref{prop:ramification-flat-family}). Statements about primary
components are made on the complex fibres; no preservation of primarity
under arbitrary nonreduced base change is asserted. The omitted locus
$\rank(K\to V)\le e-2$ is not part of this theorem. The complete
Hilbert-function-$(1,2,2)$ calculation, including its extreme-corank
locus, and the universal determinantal flag theorem remain separate
results with their original scopes.

Sections~\ref{sec:relations-ramification}--\ref{sec:intrinsic-recovery}
prove the new theorem chain. Sections~\ref{sec:v121-introduction} onward
retain the determinantal, conductor, relative and global-section
framework in full. The article is standalone; the repository designation
A2 imposes no unstated upstream or downstream hypothesis. The literature
comparison is separated from proof and is recorded in
Section~\ref{sec:v122-literature}.
